\chapter{Environnement de travail et technologies}
\label{ch:techno}

\section{Introduction}

Ce chapitre présente l'environnement matériel et logiciel utilisé pour le développement de \nomapplication, ainsi que les technologies et frameworks retenus.

\section{Environnement de travail}

\subsection{Poste de développement}

\begin{itemize}[leftmargin=*]
    \item Système d'exploitation : Windows 10/11
    \item Processeur et mémoire adaptés au développement Java et PostgreSQL
    \item IDE : Visual Studio Code / Cursor (édition) + terminal Maven
\end{itemize}

\subsection{Base de données}

\begin{itemize}[leftmargin=*]
    \item \textbf{PostgreSQL 18} installé localement (port \textbf{5433})
    \item Base : \texttt{banque\_agence}, utilisateur \texttt{banque}
    \item Script d'initialisation : \texttt{documentation/setup-postgresql.sql}
\end{itemize}

\subsection{Outils de gestion de version}

\begin{itemize}[leftmargin=*]
    \item \textbf{Git} pour le versionnement du code source
    \item Dépôt structuré par phases (commits incrémentaux)
\end{itemize}

\section{Technologies et outils utilisés}

\subsection{Back-end}

\begin{table}[H]
    \centering
    \caption{Stack back-end}
    \label{tab:stack-backend}
    \begin{tabular}{@{}lp{9cm}@{}}
        \toprule
        \textbf{Technologie} & \textbf{Rôle} \\
        \midrule
        Java 21 & Langage principal \\
        Spring Boot 3.4 & Framework applicatif (web, sécurité, JPA) \\
        Spring Security & Authentification, autorisation par rôles \\
        Spring Data JPA & Persistance objet-relationnelle (Hibernate) \\
        Flyway & Migrations versionnées de la base (V1 à V11) \\
        Bean Validation & Validation des formulaires côté serveur \\
        OpenPDF 2.0.3 & Génération des reçus PDF \\
        Maven 3.9+ & Build, dépendances, tests \\
        \bottomrule
    \end{tabular}
\end{table}

\subsection{Front-end}

\begin{table}[H]
    \centering
    \caption{Stack front-end}
    \label{tab:stack-frontend}
    \begin{tabular}{@{}lp{9cm}@{}}
        \toprule
        \textbf{Technologie} & \textbf{Rôle} \\
        \midrule
        Thymeleaf & Moteur de templates serveur (MVC) \\
        Bootstrap 5.3 & Grille responsive, composants formulaires et tableaux \\
        Bootstrap Icons 1.11 & Icônes de navigation, KPI et actions \\
        Google Fonts (DM Sans) & Typographie moderne et lisible \\
        CSS personnalisé & Charte AmanaBank (\texttt{style.css}, \texttt{landing.css}) \\
        JavaScript & En-tête fixe de la landing (\texttt{landing-header.js}) \\
        \bottomrule
    \end{tabular}
\end{table}

\subsection{Outils de conception et documentation}

\begin{itemize}[leftmargin=*]
    \item \textbf{PlantUML} : diagrammes UML (cas d'utilisation, classes, séquences, activité, MCD/MLD).
    \item \textbf{StarUML} ou éditeur texte : alternative pour la modélisation.
    \item \textbf{GanttProject} : planification des phases (optionnel).
    \item \textbf{JUnit 5 + Spring Boot Test + MockMvc} : tests d'intégration (80 tests).
\end{itemize}

\subsection{Choix technologiques}

Nous avons retenu un \textbf{monolithe Spring Boot} pour les raisons suivantes :
\begin{itemize}[leftmargin=*]
    \item Adapté à un projet PFA solo/binôme — un seul dépôt, un JAR déployable.
    \item Spring Security offre une story solide pour l'authentification bancaire.
    \item Thymeleaf évite la complexité d'un SPA React/Angular séparé.
    \item PostgreSQL garantit l'intégrité relationnelle des données financières.
    \item Flyway permet une évolution reproductible du schéma (défendable en soutenance).
\end{itemize}

\section{Conclusion}

Ce chapitre a présenté l'environnement de développement et la stack technique de \nomapplication. Le chapitre suivant détaille l'architecture logicielle et l'organisation du code source.
